% !TeX root = ../../Book1.tex
% 纲要标题：## 第20章　循环扩张、Hilbert 第 90 定理与 Kummer 理论
% 纲要位置：计划纲要.md，第 674 行。

\chapter[循环扩张、Hilbert 第 90 定理与 Kummer 理论]{循环扩张、\mbox{Hilbert 第 90 定理}与 Kummer 理论}
\label{b1:ch:20}
\BookChapterNumber{b1:ch:20}

\begin{chapterabstract}
循环扩张中的迹与范数满足特殊的相消关系。Hilbert 第 \(90\) 定理刻画其核，本章据此研究单位根条件下的根式生成与 Kummer 对应，并以 Artin--Schreier 方程处理正特征的循环扩张。
\end{chapterabstract}

\begin{learninggoals}
\begin{itemize}
\item 说明循环扩张的迹与范数公式，并区分乘法群与加法群上的两个核。
\item 掌握 Hilbert 第 90 定理的乘法形式、加法形式及各自的证明，指出可分性和嵌入独立性各在何处使用。
\item 核对 Kummer 理论的特征、单位根及有限性假设，利用配对计算扩张次数。
\item 用乘法或加法商空间中的一维子空间分类素数次循环扩张。
\end{itemize}
\end{learninggoals}

复共轭给出二次扩张 \(\mathbb C/\mathbb R\) 的一个自同构。取 \(b=2+i\)，则 \(b/\overline b=(3+4i)/5\)，这个数与其共轭的乘积为 \(1\)。若先给定一个范数为 \(1\) 的元素，是否总能反过来找到这样的 \(b\)？在循环扩张中，这个问题由 Hilbert 第 \(90\) 定理回答；它同时说明范数条件为何会出现在由幂根生成的扩张中。把乘积换成求和，又会得到一条平行而不相同的迹的结论。

\ProseParagraphBreak

本章使用\cref{b1:thm:finite-galois-correspondence}的固定域对应与\cref{b1:thm:trace-norm-embeddings}的迹、范数公式。乘法形式的 Hilbert 第 90 定理还调用\cref{b1:cor:embedding-independence}；Kummer 对应需要第一编的有限交换群结构和有限维线性代数。下一章的特征零根式可解性判据使用乘法形式、单位根与循环扩张的结果；加法形式则用于本章的 Artin--Schreier 理论。

\ProseParagraphBreak
循环扩张、二项式与 Kummer 理论，可参阅 Roman 的《Field Theory》\cite[Chapter 12; Sections 14.1--14.2, 15.1--15.3]{Roman2006FieldTheory}。若希望继续学习\subsecref{b1:subsec:2002:03}的上同调解释，可在补充群上同调之后阅读 Gille 与 Szamuely 的《Central Simple Algebras and Galois Cohomology》\cite[Example 2.3.4; Section 4.3]{GilleSzamuely2017CentralSimpleAlgebras}，其中还说明 Kummer 理论和 Artin--Schreier 理论与上同调正合列的关系。

% !TeX root = ../../Book1.tex
% 纲要标题：### 20.1 循环扩张
% 纲要位置：计划纲要.md，第 676 行。

\section{循环扩张}
\label{b1:sec:2001}

循环扩张的全部自同构都由一个自同构反复复合得到。对一个 \(n\) 次循环扩张，迹与范数分别按全部 \(n\) 个自同构求和、求积；即使不同自同构在给定元素上的取值相同，也要按出现次数计入。下面的相消计算给出迹核与范数核中的元素，下一节将证明它们恰好构成整个核。

% 纲要内容（仅作写作计划，尚非教材正文）：
%
% * 循环 Galois 群与生成自同构
% * 范数为一的元素
% * 加法与乘法形式的区别
%

\subsection{循环 Galois 群与生成自同构}
\label{b1:subsec:2001:01}

\begin{definition}\label{b1:def:cyclic-extension}
设 \(L/K\) 为有限 Galois 扩张。如果保持 \(K\) 中每个元素不变的全部 \(L\) 自同构组成的群 \(\operatorname{Gal}(L/K)\) 是循环群，则称 \(L/K\) 为\term[xunhuan-kuozhang]{循环扩张}[cyclic extension]。记 \([L:K]=n\)，并选取 \(\operatorname{Gal}(L/K)\) 的一个生成元 \(\sigma\)，于是
\[
G=\operatorname{Gal}(L/K)=\{1,\sigma,\ldots,\sigma^{n-1}\},
\qquad \sigma^n=1.
\]
\end{definition}

这里必须同时要求正规与可分。一个扩张的自同构群偶然是循环群，并不能保证它是循环扩张。例如 \(\mathbb Q(\sqrt[3]{2})/\mathbb Q\) 的自同构群只有恒等映射，但它不是正规扩张。相反，在特征不为 \(2\) 的域上，非平凡二次扩张 \(K(\sqrt d)/K\) 的两个自同构为恒等映射及 \(\sqrt d\mapsto-\sqrt d\)，故为循环扩张。

\ProseParagraphBreak
生成元通常不是唯一的。\(\sigma^r\) 生成同一个 \(n\) 阶群当且仅当 \(\gcd(r,n)=1\)。由\cref{b1:thm:finite-galois-correspondence}，每个 \(n\) 的正因子 \(d\) 对应唯一的 \(d\) 次中间扩张；具体地，\(L^{\langle\sigma^d\rangle}\) 在 \(K\) 上的次数为 \(d\)。这个次数公式使用的是子群指数，而不是子群阶。

\subsection{范数为一的元素}
\label{b1:subsec:2001:02}

设 \(L/K\) 为 \(n\) 次循环扩张，\(\operatorname{Gal}(L/K)=\langle\sigma\rangle\)。对每个 \(x\in L\)，\cref{b1:thm:trace-norm-embeddings}给出
\begin{equation}\label{b1:eq:cyclic-trace-norm}
\operatorname{Tr}_{L/K}(x)=\sum_{i=0}^{n-1}\sigma^i(x),
\qquad
N_{L/K}(x)=\prod_{i=0}^{n-1}\sigma^i(x).
\end{equation}
迹是加法群同态 \(\operatorname{Tr}_{L/K}:L\to K\)；将范数限制于非零元素，则得到群同态 \(N_{L/K}:L^\times\to K^\times\)。若 \(b\in L^\times\)，则
\[
N_{L/K}\Bigl(\frac{b}{\sigma(b)}\Bigr)
=\prod_{i=0}^{n-1}\frac{\sigma^i(b)}{\sigma^{i+1}(b)}=1.
\]
分子、分母逐项相消，末项使用 \(\sigma^n(b)=b\)。因此每个这样的比值都在范数的核中。

\ProseParagraphBreak
以 \(\mathbb C/\mathbb R\) 为例，生成自同构是复共轭，范数为 \(z\bar z\)。非零数 \(b/\bar b\) 总在单位圆上；反过来，单位圆上的每个数能否如此表示，就是乘法形式要回答的问题。这个问题只要求找出一个 \(b\)，不要求唯一：用任意非零实数乘 \(b\)，比值并不改变。

\subsection{加法形式与乘法形式}
\label{b1:subsec:2001:03}

对加法作同样计算，可得
\[
\operatorname{Tr}_{L/K}\bigl(b-\sigma(b)\bigr)
=\sum_{i=0}^{n-1}\bigl(\sigma^i(b)-\sigma^{i+1}(b)\bigr)=0.
\]
两种计算分别给出
\[
\bigl\{\,b/\sigma(b)\mid b\in L^\times\,\bigr\}\subseteq\ker N_{L/K},
\qquad
\{\,b-\sigma(b)\mid b\in L\,\}\subseteq\ker\operatorname{Tr}_{L/K}.
\]
这里的比值映射 \(b\mapsto b/\sigma(b)\) 从 \(L^\times\) 到 \(L^\times\)；必须排除 \(b=0\)，因为分母不能为零。范数映射的目标则是 \(K^\times\)。加法差映射 \(b\mapsto b-\sigma(b)\) 从 \(L\) 到 \(L\)。在特征整除 \(n\) 时，加法形式不能靠除以 \(n\) 来证明；下一节的论证依靠扩张的可分性，对这些特征仍然成立。

% !TeX root = ../../Book1.tex
% 纲要标题：### 20.2 Hilbert 第 90 定理
% 纲要位置：计划纲要.md，第 682 行。

\section{Hilbert 第 90 定理}
\label{b1:sec:2002}

上一节的相消计算给出了范数为一、迹为零的充分条件。\term[Hilbert-di90-dingli]{Hilbert 第 90 定理}[Hilbert's Theorem 90]说明这些条件也是必要的。以下固定有限循环扩张 \(L/K\)，并记 \(\operatorname{Gal}(L/K)=\langle\sigma\rangle\)、\([L:K]=n\)。

% 纲要内容（仅作写作计划，尚非教材正文）：
%
% * 循环扩张的乘法形式及证明
% * 加法形式
% * 第四卷第一群上同调的重新解释
%

\subsection{Hilbert 第 90 定理的乘法形式}
\label{b1:subsec:2002:01}

\begin{theorem}[Hilbert 第 90 定理，乘法形式]\label{b1:thm:hilbert90-multiplicative}
设 \(L/K\) 为 \(n\) 次循环扩张，并选取生成元 \(\sigma\in\operatorname{Gal}(L/K)\)。对 \(a\in L^\times\)，有 \(N_{L/K}(a)=1\) 当且仅当存在 \(b\in L^\times\) 使
\[
a=\frac{b}{\sigma(b)}.
\]
如果 \(b\) 是一个解，则全部解恰为 \(cb\)，其中 \(c\in K^\times\)。
\end{theorem}

为求解 \(a\sigma(b)=b\)，可尝试把 \(b\) 写成各 \(\sigma^i(c)\) 的加权和，并把首项权重取为 \(A_0=1\)。若使等式对每个 \(c\in L\) 成立，逐项配对得到递推 \(A_{i+1}=a\sigma(A_i)\)；递推至第 \(n\) 项时要求 \(A_n=A_0\)，恰是 \(N_{L/K}(a)=1\)。

\begin{proof}
由\eqref{b1:eq:cyclic-trace-norm}的相消计算，右边蕴含左边。现在设 \(N_{L/K}(a)=1\)，令
\[
A_0=1,\qquad A_i=a\cdot\sigma(a)\cdots\sigma^{i-1}(a)\quad(1\leq i\leq n).
\]
于是 \(A_n=1\)，且 \(a\cdot\sigma(A_i)=A_{i+1}\)。考察 \(K\) 线性映射
\[
T:L\longrightarrow L,\qquad
T(c)=\sum_{i=0}^{n-1}A_i\cdot\sigma^i(c).
\]
\cref{b1:cor:embedding-independence}说明不同嵌入 \(1,\sigma,\ldots,\sigma^{n-1}\) 作为 \(L\) 值函数线性无关。由于 \(A_0=1\)，这个映射不恒为零。选取 \(c\) 使 \(b=T(c)\neq0\)，便有
\[
a\cdot\sigma(b)=\sum_{i=0}^{n-1}A_{i+1}\cdot\sigma^{i+1}(c)
=\sum_{i=1}^{n-1}A_i\cdot\sigma^i(c)+c=b.
\]
因此 \(a=b/\sigma(b)\)。若 \(b'\) 也是解，则 \(\sigma(b'/b)=b'/b\)，故 \(b'/b\in L^{\langle\sigma\rangle}=K^\times\)；反之，乘以 \(K^\times\) 中的元素仍为解。
\end{proof}

对 \(\mathbb C/\mathbb R\)，若 \(a\bar a=1\) 且 \(a\neq-1\)，可直接取 \(b=1+a\)，因为 \(a\bar b=a+1=b\)。若 \(a=-1\)，这一选择成为零，应改取 \(b=i\)。证明中要求 \(T(c)\neq0\) 正是为了处理这种退化，不能随意把辅助元素固定为 \(1\)。

\subsection{Hilbert 第 90 定理的加法形式}
\label{b1:subsec:2002:02}

\begin{theorem}[Hilbert 第 90 定理，加法形式]\label{b1:thm:hilbert90-additive}
设 \(L/K\) 为 \(n\) 次循环扩张，并选取生成元 \(\sigma\in\operatorname{Gal}(L/K)\)。对 \(a\in L\)，有 \(\operatorname{Tr}_{L/K}(a)=0\) 当且仅当存在 \(b\in L\) 使 \(a=b-\sigma(b)\)。若 \(b\) 是一个解，则全部解组成陪集 \(b+K\)。
\end{theorem}
\begin{proof}
令 \(D:L\rightarrow L\) 为 \(K\) 线性映射 \(D(b)=b-\sigma(b)\)。其核为固定域 \(K\)，所以像的 \(K\) 维数为 \(n-1\)。另一方面，迹映射是不同嵌入之和；由\cref{b1:cor:embedding-independence}，这不是零映射。它的目标空间 \(K\) 是一维的，因此满射，核的维数也为 \(n-1\)。上一节已算出 \(\operatorname{Tr}_{L/K}\circ D=0\)，故 \(\operatorname{im}D\subseteq\ker\operatorname{Tr}_{L/K}\)。维数相等使这个包含成为等号。两个原像之差属于 \(\ker D=K\)，也就给出最后的断言。
\end{proof}

即使 \(\operatorname{char}K=p\) 且 \(p\mid n\)，迹仍然满射。此时 \(\operatorname{Tr}_{L/K}(1)=n\cdot1_K=0\) 只说明 \(1\) 落入迹的核，不能推出迹映射为零。这个区别在构造 Artin--Schreier 生成元时十分关键。

\subsection{第一群上同调的解释}
\label{b1:subsec:2002:03}

第四卷讨论群上同调时，将用\term[yici-yuquan]{一次余圈}[1-cocycle]与\term[yici-yubian]{一次余边}[1-coboundary]的语言解释这两个定理。这里仅说明有限循环 Galois 扩张中的相应方程，不要求预先学习上同调。

\ProseParagraphBreak
设 \(L/K\) 为有限 Galois 扩张，\(G=\operatorname{Gal}(L/K)\)。如果一族 \((c_g)_{g\in G}\) 中每个 \(c_g\in L^\times\)，并且满足
\[
c_{gh}=c_g\cdot g(c_h)\qquad(g,h\in G)
\]
则称它为乘法一次余圈。如果存在 \(b\in L^\times\)，使 \(c_g=b/g(b)\) 对每个 \(g\in G\) 都成立，则称这族元素为一次余边。当 \(G=\langle\sigma\rangle\)、\(|G|=n\) 时，余圈由 \(a=c_\sigma\) 决定：
\[
c_{\sigma^j}=a\cdot\sigma(a)\cdots\sigma^{j-1}(a)
\qquad(1\le j\le n).
\]
由于 \(\sigma^n=1\)，递推到 \(j=n\) 时必须有 \(N_{L/K}(a)=c_{\sigma^n}=1\)；反过来，这一条件也保证递推式与群关系相容。\cref{b1:thm:hilbert90-multiplicative}因而说明，在上述循环 Galois 群情形，每个乘法一次余圈都是一次余边。加法情形将上式改为 \(c_{gh}=c_g+g(c_h)\)，余边改为 \(b-g(b)\)；与 \(\sigma^n=1\) 相容的条件正是 \(\operatorname{Tr}_{L/K}(c_\sigma)=0\)。\cref{b1:thm:hilbert90-additive}同样说明，在此循环群情形，每个加法一次余圈都是一次余边。第四卷第十一章第 11.2 节将把一次余圈模一次余边所成的商组织成群，并重新解释这里的结论；本章的证明只使用域论。

% !TeX root = ../../Book1.tex
% 纲要标题：### 20.3 Kummer 理论
% 纲要位置：计划纲要.md，第 720 行。

\section{Kummer 理论}
\label{b1:sec:2003}

方程 \(X^n-a=0\) 看起来只涉及开方，但它的全部根还相差单位根。只有先确定基域中是否包含这些单位根，根式扩张与循环群之间的关系才足够直接。\term[Kummer-lilun]{Kummer 理论}[Kummer theory]研究这类根式与交换 Galois 群的联系。本节固定代数闭包 \(\overline K\)，所有扩张都取在其中；除另有说明外，\(n\geq1\)、\(\operatorname{char}K\nmid n\)，并假设 \(K\) 含有全部 \(n\) 次单位根。

% 纲要内容（仅作写作计划，尚非教材正文）：
%
% * 特征不整除 n 且基域含所需单位根的假设
% * 先以 Hilbert 第 90 定理证明循环扩张的根式生成，作为第21.2节的直接准备
% * 单根式扩张次数等于被开方数的商类阶；根式参数选择与 Kummer 配对的良定义
% * 有限 Kummer 对应的双侧单射、阶比较与同时特征基重建；素数次扩张由参数直线分类
% * 在本节定义 \(\mu_n\) 及本原单位根；循环扩张的根式生成不依赖一般 Kummer 对应，也不调用后续分圆多项式不可约性
%

\subsection{特征与单位根假设}
\label{b1:subsec:2003:01}

\begin{definition}\label{b1:def:roots-of-unity}
设 \(K\) 为域，\(\overline K\) 为 \(K\) 的代数闭包，\(n\geq1\) 为整数。如果 \(\zeta\in\overline K^\times\) 满足 \(\zeta^n=1\)，则称 \(\zeta\) 为 \(n\) 次\term[danwei-gen]{单位根}[root of unity]。全体 \(n\) 次单位根组成乘法群
\[
\mu_n\coloneqq\{\,\zeta\in\overline K^\times\mid\zeta^n=1\,\}.
\]
如果 \(\zeta\in\mu_n\) 的乘法阶恰为 \(n\)，则称 \(\zeta\) 为\term[benyuan-danwei-gen]{本原单位根}[primitive root of unity]。
\end{definition}
\symidx{\(\mu_n\)：全体 \(n\) 次单位根组成的群}

在 \(\operatorname{char}K\nmid n\) 的假设下，\(X^n-1\) 与导数 \(nX^{n-1}\) 没有公共根，所以 \(\mu_n\) 有 \(n\) 个元素。\cref{b1:lem:finite-multiplicative-cyclic}说明域的乘法群的有限子群是循环群；\(\mu_n\) 因而是阶为 \(n\) 的循环群，它的任一生成元 \(\zeta\) 都有阶 \(n\)，即为本原 \(n\) 次单位根，且 \(\mu_n=\langle\zeta\rangle\)。假设 \(\mu_n\subseteq K\) 就等价于 \(K\) 含有这样一个 \(\zeta\)。

\ProseParagraphBreak
对 \(a\in K^\times\)，特征条件使 \((X^n-a)'=nX^{n-1}\) 与 \(X^n-a\) 没有公共根，保证多项式可分；若去掉它，在特征 \(p\) 中就有 \(X^p-a=(X-\alpha)^p\)，其中 \(\alpha^p=a\)，只有一个重根，并且 \(X^p-1=(X-1)^p\) 表明不存在阶为 \(p\) 的单位根。单位根条件则保证：已有一个根 \(\alpha\) 后，其余根 \(\zeta\alpha\)（\(\zeta\in\mu_n\)）全部位于 \(K(\alpha)\) 中。若去掉它，\(\mathbb Q(\sqrt[3]{2})/\mathbb Q\) 就是可分而不正规的根式扩张。这两个假设分别排除了重根与遗漏其他共轭根的障碍。

\subsection{循环扩张与根式生成}
\label{b1:subsec:2003:04}

\begin{proposition}\label{b1:prop:single-radical-cyclic}
设 \(K\) 为域，\(n\geq1\) 为整数，\(\operatorname{char}K\nmid n\)，且 \(K\) 包含全部 \(n\) 次单位根。固定 \(K\) 的代数闭包 \(\overline K\)，取 \(a\in K^\times\) 及满足 \(\alpha^n=a\) 的 \(\alpha\in\overline K\)。则 \(K(\alpha)/K\) 是循环 Galois 扩张，其次数整除 \(n\)。映射
\[
\operatorname{Gal}\bigl(K(\alpha)/K\bigr)\longrightarrow\mu_n,
\qquad\sigma\longmapsto\frac{\sigma(\alpha)}{\alpha}
\]
是单同态。
\end{proposition}
\begin{proof}
取 \(\mu_n\) 的生成元 \(\zeta\)。\(X^n-a\) 的根恰为 \(\zeta^j\alpha\)，全在 \(K(\alpha)\) 中；因 \(a\neq0\) 且 \(n\neq0\) 在 \(K\) 中成立，它没有重根。所以 \(K(\alpha)/K\) 是一个可分多项式的分裂域，为有限 Galois 扩张。对其任一自同构 \(\sigma\)，
\[
\sigma(\alpha)^n=\sigma(\alpha^n)=\sigma(a)=a=\alpha^n,
\]
故 \((\sigma(\alpha)/\alpha)^n=1\)，比值属于 \(\mu_n\)。又因 \(\mu_n\subseteq K\)，每个自同构都固定这些比值，所以
\[
\frac{\sigma\tau(\alpha)}{\alpha}
=\sigma\Bigl(\frac{\tau(\alpha)}{\alpha}\Bigr)\frac{\sigma(\alpha)}{\alpha}
=\frac{\tau(\alpha)}{\alpha}\frac{\sigma(\alpha)}{\alpha}.
\]
目标群交换，故上式就是普通群同态的乘法公式。比值为 \(1\) 时，自同构固定生成元 \(\alpha\)，因而为恒等映射，这证明了单射性。Galois 群由此同构于 \(\mu_n\) 的一个子群，所以为循环群，群阶整除 \(|\mu_n|=n\)。由于扩张已经证明为 Galois，\([K(\alpha):K]=|\operatorname{Gal}(K(\alpha)/K)|\)，扩张次数也整除 \(n\)。
\end{proof}

这个同态计算再次用到了单位根条件。在一个更大的分裂域中，若自同构不固定单位根，则 \(\sigma(\tau(\alpha)/\alpha)\) 不能直接换成 \(\tau(\alpha)/\alpha\)，仍须保留自同构对该比值的作用；那正是上一节一次余圈的乘法公式。

\ProseParagraphBreak
根式扩张的次数还可以从被开方数在乘法商群中的阶直接读出。

\begin{corollary}\label{b1:cor:kummer-radical-degree}
设 \(K\) 为域，\(n\geq1\) 为整数，\(\operatorname{char}K\nmid n\)，且 \(K\) 包含全部 \(n\) 次单位根。固定 \(K\) 的代数闭包，取 \(a\in K^\times\)，并在该闭包中取满足 \(\alpha^n=a\) 的根 \(\alpha\)。记 \(K^{\times n}=\{\,c^n\mid c\in K^\times\,\}\)。则
\[
[K(\alpha):K]=\operatorname{ord}_{K^\times/K^{\times n}}([a]).
\]
\end{corollary}
\begin{proof}
记 \(L=K(\alpha)\)、\(d=[L:K]\)，并令 \(e\) 为 \([a]\) 的阶。因为 \([a]^n=[1]\)，\(e\) 是整除 \(n\) 的正整数。由\cref{b1:prop:single-radical-cyclic}，\(L/K\) 循环 Galois，其 Galois 群通过比值映射嵌入 \(\mu_n\)，像群的阶为 \(d\)。因此每个比值 \(\sigma(\alpha)/\alpha\) 的 \(d\) 次幂为一，故 \(\sigma(\alpha^d)=\alpha^d\) 对所有自同构成立。不动域为 \(K\)，所以 \(\alpha^d\in K^\times\)，进而 \(a^d=(\alpha^d)^n\in K^{\times n}\)，即 \(e\mid d\)。

反过来，\(a^e=c^n\) 对某个 \(c\in K^\times\) 成立，故 \((\alpha^e/c)^n=1\)。由 \(\mu_n\subseteq K\) 得 \(\alpha^e/c\in K\)，所以 \(\alpha^e\in K\)。于是 \(\alpha\) 满足 \(K\) 上的 \(e\) 次多项式 \(X^e-\alpha^e\)，其最小多项式次数 \(d\) 不超过 \(e\)。结合 \(e\mid d\) 与 \(d\le e\)，得到 \(d=e\)。
\end{proof}

反过来，循环扩张也可以用一个根式生成。

\begin{theorem}\label{b1:thm:kummer-cyclic}
设 \(L/K\) 是次数为 \(d\) 的循环扩张。如果 \(\operatorname{char}K\nmid d\)，并且 \(K\) 包含全部 \(d\) 次单位根 \(\mu_d\)，则存在 \(\alpha\in L^\times\) 使
\[
L=K(\alpha),\qquad\alpha^d\in K.
\]
\end{theorem}

为寻找这样的 \(\alpha\)，选 Galois 群的生成元 \(\sigma\) 和本原 \(d\) 次单位根 \(\zeta\)，要求 \(\sigma(\alpha)=\zeta\alpha\)，就是在线性代数意义下寻找 \(\sigma\) 的一个特征向量，特征值为 \(\zeta\)。这使 \(\sigma(\alpha^d)=\zeta^d\alpha^d=\alpha^d\)，又把构造化为 \(\alpha/\sigma(\alpha)=\zeta^{-1}\)，恰可调用乘法形式的 Hilbert 第 90 定理。

\begin{proof}
选取 \(\operatorname{Gal}(L/K)\) 的生成元 \(\sigma\)，以及本原 \(d\) 次单位根 \(\zeta\in K\)。因为 \(\zeta\) 属于基域，全部自同构都固定它，所以
\[
N_{L/K}(\zeta^{-1})=\prod_{j=0}^{d-1}\sigma^j(\zeta^{-1})
=\zeta^{-d}=1.
\]
\cref{b1:thm:hilbert90-multiplicative}给出非零 \(\alpha\) 满足 \(\alpha/\sigma(\alpha)=\zeta^{-1}\)，即 \(\sigma(\alpha)=\zeta\alpha\)。所以 \(\alpha^d\) 被 \(\sigma\) 固定，属于 \(K\)。又有 \(\sigma^j(\alpha)=\zeta^j\alpha\) 对 \(0\le j<d\) 成立。\(\sigma^j\) 固定最小多项式 \(m_{\alpha,K}\) 的系数，故这些像全部是它的根；因为 \(\zeta\) 的阶为 \(d\) 且 \(\alpha\ne0\)，它们两两不同。因此 \(\deg m_{\alpha,K}\ge d\)。另一方面，\(K(\alpha)\subseteq L\) 给出 \(\deg m_{\alpha,K}=[K(\alpha):K]\le[L:K]=d\)，故次数相等，\(K(\alpha)=L\)。
\end{proof}

\cref{b1:thm:kummer-cyclic}将用于\secref{b1:sec:2102}的特征零根式可解性判据。多个根式之间的关系还可由有限交换 Galois 群统一描述，下面据此建立 Kummer 配对。

\subsection{根式扩张与 Kummer 配对}
\label{b1:subsec:2003:02}

一个根式的自同构由比值 \(\sigma(\alpha)/\alpha\) 决定。加入多个根式时，我们希望同时记录这些比值，并判断哪些根式增加了扩张次数。若 \(a'=ab^n\)、\(b\in K^\times\)，而 \(\alpha^n=a\)，则 \((b\alpha)^n=a'\)，且 \(K(b\alpha)=K(\alpha)\)。因此应先把被开方数放入商群 \(K^\times/K^{\times n}\)，再用同一自同构在各个根上的比值来比较它们。

\ProseParagraphBreak
为同时处理若干根式，记 \(K^{\times n}=\{\,b^n\mid b\in K^\times\,\}\)。商群 \(K^\times/K^{\times n}\) 中每个元素的阶都整除 \(n\)。取其有限子群 \(B\)，选择生成元的代表 \(a_1,\ldots,a_r\)，并令
\[
L_B=K(\alpha_1,\ldots,\alpha_r),\qquad\alpha_i^n=a_i.
\]
换一个根只会乘以 \(\mu_n\subseteq K\) 中的元素；换一个代表只会再乘以 \(K^\times\) 中的元素。若 \([a]\in B\)，可写 \([a]=[a_1]^{e_1}\cdots[a_r]^{e_r}\)，其中 \(e_i\in\mathbb Z\)。这意味着对某个 \(c\in K^\times\) 有
\[
a=c^n a_1^{e_1}\cdots a_r^{e_r},\qquad
\bigl(c\alpha_1^{e_1}\cdots\alpha_r^{e_r}\bigr)^n=a.
\]
所有 \(\alpha_i\) 均非零，所以负指数也有意义；\(a\) 的任一 \(n\) 次根只与所列的根相差 \(K\) 中的单位根，故都属于 \(L_B\)。换一组生成元时，新生成元代表的根因而全部属于原来的 \(L_B\)，反过来原生成元的根也属于新域，所以两域相等。这说明 \(L_B\) 不依赖代表、根或生成列表的选择。它也是多项式 \(X^n-a_1,\ldots,X^n-a_r\) 的共同分裂域；这些多项式都可分，故 \(L_B/K\) 是有限 Galois 扩张。记 \(G_B=\operatorname{Gal}(L_B/K)\)。

\begin{definition}\label{b1:def:kummer-pairing}
设 \(K\) 为域，\(n\geq1\) 且 \(\operatorname{char}K\nmid n\)，在 \(K\) 的固定代数闭包 \(\overline K\) 中记 \(\mu_n=\{\zeta\in\overline K^\times\mid\zeta^n=1\}\)，并假设 \(\mu_n\subseteq K\)。取乘法商群 \(K^\times/K^{\times n}\) 的有限子群 \(B\)，其中 \(K^{\times n}=\{b^n\mid b\in K^\times\}\)；令 \(L_B\) 为在 \(\overline K\) 中给 \(B\) 的生成元代表添入 \(n\) 次根所得的域，并记 \(G_B=\operatorname{Gal}(L_B/K)\)。对 \(\sigma\in G_B\) 与 \(\xi=[a]\in B\)，取 \(a\in K^\times\) 及满足 \(\alpha^n=a\) 的 \(\alpha\in L_B\)，定义
\[
\langle\sigma,\xi\rangle\coloneqq\frac{\sigma(\alpha)}{\alpha},
\qquad \alpha^n=a.
\]
该值与代表 \(a\) 及根 \(\alpha\) 的选择无关。由此得到的映射 \(G_B\times B\rightarrow\mu_n\) 称为\term[Kummer-peidui]{Kummer 配对}[Kummer pairing]。
\end{definition}

配对的值与根和代表的选择无关。若把 \(\alpha\) 换成 \(\alpha'=\zeta\alpha\)，其中 \(\zeta\in\mu_n\subseteq K\)，则
\[
\frac{\sigma(\alpha')}{\alpha'}=
\frac{\zeta\sigma(\alpha)}{\zeta\alpha}=
\frac{\sigma(\alpha)}{\alpha}.
\]
若把代表换成 \(a'=ab^n\)，其中 \(b\in K^\times\)，可取根 \(b\alpha\)，同样有 \(\sigma(b\alpha)/(b\alpha)=\sigma(\alpha)/\alpha\)。该代表的其他根再由前一种计算处理，故两个选择都不影响配对值。

\ProseParagraphBreak
根的乘积是乘积的根，所以配对在第二个变量上保持乘法；第一个变量上的计算仍须使用 \(\sigma\) 固定 \(\mu_n\) 这一事实，正如\cref{b1:prop:single-radical-cyclic}证明中的比值公式。于是对 \(\sigma,\tau\in G_B\)、\(\xi,\eta\in B\)，有
\[
\begin{aligned}
\langle\sigma\tau,\xi\rangle&=
\langle\sigma,\xi\rangle\langle\tau,\xi\rangle,\\
\langle\sigma,\xi\eta\rangle&=
\langle\sigma,\xi\rangle\langle\sigma,\eta\rangle.
\end{aligned}
\]
因此固定任意一个变量，就得到另一个变量上的群同态。

\subsection{有限 Kummer 对应}
\label{b1:subsec:2003:03}

在 \(\operatorname{char}K\nmid n\) 的假设下，\(\mu_n\) 是阶为 \(n\) 的循环群。设 \(A\) 为有限交换群。对群同态 \(\chi,\psi:A\rightarrow\mu_n\)，\(\operatorname{Hom}(A,\mu_n)\) 的运算是逐点相乘：\((\chi\psi)(x)=\chi(x)\psi(x)\)，单位元是常值 \(1\) 的同态。\(A\) 的指数指使每个元素的该次幂等于 \(1\) 的最小正整数。若其指数整除 \(n\)，则
\begin{equation}\label{b1:eq:finite-character-count}
\bigl|\operatorname{Hom}(A,\mu_n)\bigr|=|A|.
\end{equation}
这里应用\cref{b1:thm:fga-structure}，把有限交换群 \(A\) 写成若干循环群的直积；因 \(A\) 的指数整除 \(n\)，每个循环因子的阶 \(d\) 也整除 \(n\)。从 \(C_d\) 到 \(\mu_n\) 的同态由生成元的像决定，这个像必须满足 \(\zeta^d=1\)。在阶为 \(n\) 的循环群中，满足这个条件的元素有 \(\gcd(d,n)\) 个；只有 \(d\mid n\) 时才恰有 \(d\) 个。当前每个循环因子均满足此条件，各因子的选择独立，相乘得到所需等式。这说明指数条件不能省去；单凭 \(A\) 有限交换，并不能保证有 \(|A|\) 个这样的同态。

\begin{theorem}[Kummer 对应]\label{b1:thm:kummer-correspondence}
设 \(K\) 为域，\(n\geq1\) 为整数，\(\operatorname{char}K\nmid n\)，并且 \(K\) 包含全部 \(n\) 次单位根 \(\mu_n\)。固定 \(K\) 的代数闭包 \(\overline K\)，记 \(K^{\times n}=\{c^n\mid c\in K^\times\}\)。以下两类对象一一对应：
\begin{enumerate}
\item \(K^\times/K^{\times n}\) 的有限子群 \(B\)；
\item \(\overline K\) 中的有限 Galois 扩张 \(L/K\)，其 Galois 群为指数整除 \(n\) 的交换群。
\end{enumerate}
对第一类中的有限子群 \(B\)，令 \(L_B\subseteq\overline K\) 为给 \(B\) 中各元素的代表添入 \(n\) 次根所生成的域。正向对应为 \(B\mapsto L_B\)，反向对应为
\[
B_L=\left(K^\times\cap L^{\times n}\right)/K^{\times n}.
\]
这里 \(L^{\times n}=\{\,\beta^n\mid\beta\in L^\times\,\}\)，所以分子收集 \(K\) 中那些在 \(L\) 内已经具有 \(n\) 次根的非零元素，再模掉它们在 \(K\) 中本来就是 \(n\) 次幂的平凡类。
Kummer 配对给出同构
\[
\operatorname{Gal}(L_B/K)\cong\operatorname{Hom}(B,\mu_n),
\qquad B\cong\operatorname{Hom}\bigl(\operatorname{Gal}(L_B/K),\mu_n\bigr),
\]
特别地，\([L_B:K]=|B|\)。子群包含对应于域包含。
\end{theorem}
\begin{proof}
从有限 \(B\) 出发，配对给出群同态
\[
G_B\longrightarrow\operatorname{Hom}(B,\mu_n),\qquad
\sigma\longmapsto\bigl(\xi\mapsto\langle\sigma,\xi\rangle\bigr).
\]
若 \(\sigma\) 在其核中，则全部选定根的比值 \(\sigma(\alpha_i)/\alpha_i\) 均为一。因此 \(\sigma\) 固定生成域的每个根 \(\alpha_i\)，也固定由它们和 \(K\) 生成的整个 \(L_B\)，故为恒等。这个映射因而单射。目标是逐点相乘的交换群，每个同态的 \(n\) 次幂都是常值一的同态，所以这个单射还说明 \(G_B\) 交换且指数整除 \(n\)。

在另一侧，配对给出 \(B\rightarrow\operatorname{Hom}(G_B,\mu_n)\)。若 \([a]\) 对每个 \(\sigma\in G_B\) 的配对值都是一，则其根 \(\alpha\) 被整个群固定。不动域为 \(K\)，所以 \(\alpha\in K\)，\(a=\alpha^n\) 在 \(K\) 中本来就是 \(n\) 次幂，即 \([a]=[1]\)。这也得到一个单射：任何非平凡被开方数类都能被某个自同构通过配对区别出来。

有限交换群 \(B\) 的指数整除 \(n\)，所以\eqref{b1:eq:finite-character-count}给出 \(|\operatorname{Hom}(B,\mu_n)|=|B|\)，第一个单射推出 \(|G_B|\le|B|\)。刚才已从这个单射得到 \(G_B\) 的交换性与指数条件，因而同一个计数公式也适用于 \(G_B\)，得到 \(|\operatorname{Hom}(G_B,\mu_n)|=|G_B|\)。第二个单射于是推出 \(|B|\le|G_B|\)。两边的有限阶相等，两个单射都为满射，故都是同构。又因 \(L_B/K\) 为 Galois，\([L_B:K]=|G_B|=|B|\)。

现在从第二类扩张 \(L/K\) 出发，记群为 \(G\)。要将它重建为一个根式扩张，只需找到一组 \(K\) 基 \(\beta_i\)，使每个 \(\beta_i^n\) 都在 \(K\) 中。把各 \(\sigma\in G\) 视为 \(K\) 向量空间 \(L\) 上的线性算子，寻找它们的共同特征向量，就能使这些 \(n\) 次幂被整个群固定。由于群的指数整除 \(n\)，各算子满足 \(\sigma^n=1\)；\(X^n-1\) 在 \(K\) 中分裂且无重根，可以用插值多项式把向量分解到其特征空间中。对每个 \(\zeta\in\mu_n\)，令
\[
P_\zeta(X)=\prod_{\substack{\eta\in\mu_n\\\eta\neq\zeta}}
\frac{X-\eta}{\zeta-\eta}.
\]
这就是\cref{b1:thm:lagrange-interpolation}在互异插值点 \(\mu_n\) 上的基多项式，满足 \(P_\zeta(\eta)=1\) 当 \(\eta=\zeta\)，否则为零。因而 \(\sum_\zeta P_\zeta(X)-1\) 是次数小于 \(n\) 却有 \(n\) 个不同根的多项式，只能为零。又有
\[
(X-\zeta)P_\zeta(X)=
\frac{X^n-1}{\prod_{\eta\in\mu_n,\,\eta\ne\zeta}(\zeta-\eta)},
\]
因为分子恰是全部 \(n\) 个一次因子的乘积，分母是非零的基域常数。代入算子 \(\sigma\)，利用 \(\sigma^n=1\)，得到 \((\sigma-\zeta I)P_\zeta(\sigma)=0\)。所以
\[
v=\sum_{\zeta\in\mu_n}P_\zeta(\sigma)v
\]
的每一项都属于 \(\sigma\) 的 \(\zeta\) 特征空间。\(P_\zeta(\sigma)\) 在该空间上为恒等，在其他特征空间上为零；若这些空间中向量之和为零，分别作用各 \(P_\zeta(\sigma)\) 就得到每项都为零。因此这个分解是直和，\(\sigma\) 可对角化。

要得到整个 \(G\) 的同时特征基，还需要群交换。若 \(\sigma(v)=\zeta v\)，则对任意 \(\tau\in G\)，
\[
\sigma(\tau(v))=\tau(\sigma(v))=\tau(\zeta v)=\zeta\tau(v),
\]
因为两算子交换且 \(\zeta\in K\)。因此一个算子的每个特征空间都被其余算子保持，可以在其中对下一算子再作上述直和分解。群 \(G\) 有限，依次处理其全部元素后，便得到 \(L\) 的公共特征空间直和分解；在每个非零公共特征空间中选一组基，合在一起得到同时特征向量基 \(\beta_1,\ldots,\beta_m\)。这里指数条件用于单个算子的分解，交换性用于使这些分解相容。

写 \(\sigma(\beta_i)=\lambda_{\sigma,i}\beta_i\)，其中 \(\lambda_{\sigma,i}\in K\)。因为 \(\beta_i\ne0\)，
\[
\beta_i=\sigma^n(\beta_i)=\lambda_{\sigma,i}^n\beta_i
\]
给出 \(\lambda_{\sigma,i}^n=1\)，即 \(\lambda_{\sigma,i}=\sigma(\beta_i)/\beta_i\in\mu_n\)。于是 \(\sigma(\beta_i^n)=\beta_i^n\) 对所有 \(\sigma\in G\) 成立，故 \(\beta_i^n\in K^\times\)。这些 \(n\) 次幂的类生成一个有限子群 \(B_0\)：各生成元的幂指数都可约为 \(0,\ldots,n-1\)，所以其元素数至多为 \(n^m\)。构造 \(L_{B_0}\) 时可取 \(\beta_i\) 作为这些类的根，因此 \(L_{B_0}\subseteq L\)。另一方面，这个子域含有 \(K\) 和全部 \(\beta_i\)，故含有它们的所有 \(K\) 线性组合，即整个 \(L\)。两边包含给出 \(L_{B_0}=L\)。

\(B_0\) 是利用一次基的选择构造出来的参数群；还须确认它已经包含 \(L\) 中所有被开方数类，即等于由扩张本身定义的 \(B_L\)。对 \(B_L\) 中任意元素，在 \(L\) 中选取其代表的 \(n\) 次根。以同一比值公式得到群同态 \(B_L\rightarrow\operatorname{Hom}(G,\mu_n)\)；若某个类与所有 \(\sigma\in G\) 的配对值均为一，其根便属于不动域 \(K\)，所以该映射单射。目标是有限群，因而 \(B_L\) 有限。由构造可知 \(B_0\subseteq B_L\)；而 \(B_L\) 中各类在 \(L\) 中有根，另选的根只差 \(K\) 中的单位根，因此 \(L_{B_L}\subseteq L\)。于是
\[
L=L_{B_0}\subseteq L_{B_L}\subseteq L.
\]
故 \(L_{B_L}=L\)。对两个有限子群分别应用已证的次数公式，得到 \(|B_L|=[L:K]=|B_0|\)。因为 \(B_0\subseteq B_L\)，有限群的阶相等迫使 \(B_L=B_0\)，所以选取的这组基虽然不唯一，所得到的参数群却是扩张内在确定的全部参数群。

反过来，从任意有限 \(B\) 出发，由各代表在 \(L_B\) 中具有根可知 \(B\subseteq B_{L_B}\)。刚才对第二类扩张的证明应用于 \(L_B/K\)，给出 \(B_{L_B}\) 有限且 \(L_{B_{L_B}}=L_B\)。因此 \(|B_{L_B}|=[L_B:K]=|B|\)，由包含与阶相等得 \(B_{L_B}=B\)。这证明了两个构造互逆。若 \(B\subseteq B'\)，则其代表的根也包含于 \(L_{B'}\)，从而 \(L_B\subseteq L_{B'}\)；反之，若 \(L_B\subseteq L_{B'}\)，一个在前者中已有 \(n\) 次根的基域元素在后者中也有根，故 \(B=B_{L_B}\subseteq B_{L_{B'}}=B'\)。
\end{proof}

这里参数子群 \(B\) 增大，所要添入的根也增多，所以 Kummer 对应中的子群包含与域包含同向。上一章的 Galois 对应则反向：自同构子群越大，共同不动元需要满足的条件越多，不动域越小。两种对应中的子群分别记录被开方数类与自同构，含义不同。

\ProseParagraphBreak
定理要求扩张有限，且 Galois 群的指数整除给定的 \(n\)。与之对应的 \(B\) 是有限子群，但不要求它在整个 \(K^\times/K^{\times n}\) 中具有有限指数。例如取 \(K=\mathbb Q\)、\(n=2\)，则 \(\mathbb Q^\times/\mathbb Q^{\times2}\) 是无限群，但 \(B=\{[1],[2],[3],[6]\}\) 有四个元素，给出 \(L_B=\mathbb Q(\sqrt2,\sqrt3)\)。独立性可由 \(2,3,6\) 都不是有理数的平方核验。

\ProseParagraphBreak
令 \(s\) 只改变 \(\sqrt2\) 的符号，\(t\) 只改变 \(\sqrt3\) 的符号。此时 \(G_B=\{1,s,t,st\}\)、\(\mu_2=\{1,-1\}\)，配对的全部取值为
\[
\begin{array}{c|rrrr}
\langle\sigma,\xi\rangle &[1]&[2]&[3]&[6]\\ \hline
1  & 1& 1& 1& 1\\
s  & 1&-1& 1&-1\\
t  & 1& 1&-1&-1\\
st & 1&-1&-1& 1
\end{array}
\]
每一行是 \(B\rightarrow\mu_2\) 的同态，每一列是 \(G_B\rightarrow\mu_2\) 的同态；各行两两不同，各列也两两不同，具体显示了配对的两侧非退化性。特别地，\([6]\) 的核为 \(\{1,st\}\)，所以 \(\mathbb Q(\sqrt6)=L_B^{\langle st\rangle}\)。这里 \(\{[1],[6]\}\subseteq B\) 给出子域包含 \(\mathbb Q(\sqrt6)\subseteq L_B\)，而子域与其固定子群之间的包含方向相反。

\ProseParagraphBreak
设 \(K\) 为域，\(\ell\) 为素数。乘法商群 \(V=K^\times/K^{\times\ell}\) 的每个元素的阶都整除 \(\ell\)，因而可以把 \(V\) 视为 \(\mathbb F_\ell\) 向量空间；\(V\) 也可能是平凡群。若用加法记号，其运算为
\[
[a]+[b]=[ab],\qquad \bar r\,[a]=[a^r],\qquad 0=[1]
\quad(a,b\in K^\times,\ r\in\mathbb Z).
\]
其中 \(\bar r\) 是整数 \(r\) 在 \(\mathbb F_\ell\) 中的像。因为 \([a]^\ell=[1]\)，若 \(r'\equiv r\pmod\ell\)，则 \([a^{r'}]=[a^r]\)，所以标量乘法只依赖 \(\bar r\)，不依赖其整数代表。分配律与标量乘法的结合律来自交换群的幂运算规则，例如 \((ab)^r=a^rb^r\)、\(a^{r+s}=a^ra^s\)。这里的非零向量是 \([a]\ne[1]\)，即 \(a\in K^\times\setminus K^{\times\ell}\)；只写 \(a\ne0\) 并不足以保证这一点。


\begin{corollary}\label{b1:cor:kummer-prime-classification}
设 \(K\) 为域，\(\ell\) 为不同于 \(\operatorname{char}K\) 的素数，且 \(K\) 包含全部 \(\ell\) 次单位根 \(\mu_\ell\)。在 \(K\) 的一个固定代数闭包中，\(K\) 的 \(\ell\) 次循环扩张由 \(K^\times/K^{\times\ell}\) 的一维 \(\mathbb F_\ell\) 子空间分类。
\end{corollary}
\begin{proof}
素数情形应用\cref{b1:thm:kummer-correspondence}：次数为 \(\ell\) 的域对应阶为 \(\ell\) 的子群。上述乘法商群的指数整除 \(\ell\)，它作为 \(\mathbb F_\ell\) 向量空间时，阶为 \(\ell\) 的子群正是一维子空间。
\end{proof}

因此，若 \(a\in K^\times\setminus K^{\times\ell}\)，同一直线上的 \(\ell-1\) 个非零向量 \([a]^r\)（\(1\le r<\ell\)）都生成同一个参数子群 \(\langle[a]\rangle\)，给出同一个非平凡 \(\ell\) 次扩张。具体地，若 \(\alpha^\ell=a\)、\(b\in K^\times\)，则 \(b\alpha^r\) 的 \(\ell\) 次幂为 \(a^r b^\ell\)。由于 \(r\) 与 \(\ell\) 互素，取整数 \(u,v\) 使 \(ru+\ell v=1\)，就有
\[
\alpha=(\alpha^r)^u a^v\in K(b\alpha^r).
\]
反向包含显然，所以 \(K(b\alpha^r)=K(\alpha)\)。以每个非零类分别标记扩张，会将这一条直线重复计算 \(\ell-1\) 次；一维子空间恰好把这些生成元变化归为同一个参数。这个区别与下一节的加法分类相对应。

% !TeX root = ../../Book1.tex
% 纲要标题：### 20.4 Artin–Schreier 理论
% 纲要位置：计划纲要.md，第 728 行。

\section{Artin–Schreier 理论}
\label{b1:sec:2004}

在特征 \(p\) 中，\(p\) 次开方通常产生纯不可分扩张，不能代替 \(p\) 次循环 Galois 扩张。一个导数恒为 \(-1\) 的多项式提供了正确的可分模型。以下 \(K\) 始终为特征 \(p>0\) 的域，扩张取在一个固定代数闭包中。

% 纲要内容（仅作写作计划，尚非教材正文）：
%
% * 特征 p 中的方程 X^p−X−a；素域线性、根的平移与不可约二分
% * p 次循环扩张的构造与参数直线分类；用有理函数次数解释无穷远极点
% * 与 Kummer 理论的平行及差别；区分纯不可分开方与可分循环步骤
%
% ---
%

\subsection{Artin–Schreier 方程}
\label{b1:subsec:2004:01}

\begin{definition}\label{b1:def:artin-schreier-map}
设 \(K\) 为特征 \(p>0\) 的域。映射
\[
\wp:K\longrightarrow K,\qquad b\longmapsto b^p-b
\]
称为\term[Artin-Schreier-yingshe]{Artin--Schreier 映射}[Artin--Schreier map]。给定 \(a\in K\)，把方程 \(X^p-X-a=0\) 称为相应的\term[Artin-Schreier-fangcheng]{Artin--Schreier 方程}[Artin--Schreier equation]。
\end{definition}
\symidx{\(\wp(b)=b^p-b\)：Artin--Schreier 映射}

在特征 \(p\) 中，二项式公式给出 \((u+v)^p=u^p+v^p\)。又因 \(c^p=c\) 对每个 \(c\in\mathbb F_p\) 成立，有
\[
\wp(u+v)=\wp(u)+\wp(v),\qquad
\wp(cu)=c^pu^p-cu=c\wp(u).
\]
这证明了 \(\wp\) 的 \(\mathbb F_p\) 线性性。

\ProseParagraphBreak
素域 \(\mathbb F_p\) 的 \(p\) 个元素都是 \(X^p-X\) 的根；这个多项式的次数为 \(p\)，在任何扩域中都没有其他根。因此 \(\wp(b)=0\) 当且仅当 \(b\in\mathbb F_p\)，即 \(\ker\wp=\mathbb F_p\)。若 \(K\neq\mathbb F_p\)，取 \(d\in K\setminus\mathbb F_p\)，则 \(\wp(d\cdot1)=d^p-d\neq0=d\wp(1)\)，所以它并不是 \(K\) 线性映射。像 \(\wp(K)\) 是 \(K\) 的加法群中的 \(\mathbb F_p\) 线性子空间，所以 \(K/\wp(K)\) 是 \(\mathbb F_p\) 向量空间。后文的一维子空间正是按这个线性结构理解的；这里没有对乘法取商，也不把它看成商域。

\begin{proposition}\label{b1:prop:artin-schreier-polynomial}
设 \(K\) 为特征 \(p>0\) 的域，\(a\in K\)，记 \(\wp(K)=\{b^p-b\mid b\in K\}\)。则 \(X^p-X-a\) 要么在 \(K\) 上完全分裂，要么在 \(K[X]\) 中不可约。前者当且仅当 \(a\in\wp(K)\)；在后一种情形，该多项式的任一根 \(\alpha\) 都生成 \(p\) 次循环扩张 \(K(\alpha)/K\)，其全部 \(K\) 自同构为 \(\alpha\mapsto\alpha+c\)，其中 \(c\in\mathbb F_p\)。
\end{proposition}
\begin{proof}
若 \(\alpha,\beta\) 都是 \(X^p-X-a\) 的根，则
\[
(\beta-\alpha)^p-(\beta-\alpha)
=(\beta^p-\beta)-(\alpha^p-\alpha)=0.
\]
由 \(X^p-X\) 的根的描述，\(\beta-\alpha\in\mathbb F_p\)。反过来，每个 \(\alpha+c\)、\(c\in\mathbb F_p\)，都满足原方程。因此全部根恰为这 \(p\) 个不同的元素。它们都在 \(L=K(\alpha)\) 中，所以 \(L\) 是该多项式的分裂域；导数为 \(-1\)，又保证它没有重根。由可分多项式的分裂域判据，\(L/K\) 是有限 Galois 扩张，故\cref{b1:prop:galois-automorphism-count}给出 \(|\operatorname{Gal}(L/K)|=[L:K]\)。

对 \(\sigma\in\operatorname{Gal}(L/K)\)，令 \(c_\sigma=\sigma(\alpha)-\alpha\)。因为 \(\sigma(\alpha)\) 也是原方程的根，所以 \(c_\sigma\in\mathbb F_p\)。每个 \(K\) 自同构都固定素域，故
\[
\begin{aligned}
c_{\sigma\tau}
&=\sigma\tau(\alpha)-\alpha\\
&=\sigma\bigl(\tau(\alpha)-\alpha\bigr)+\sigma(\alpha)-\alpha\\
&=c_\tau+c_\sigma.
\end{aligned}
\]
这说明 \(\sigma\mapsto c_\sigma\) 是到 \((\mathbb F_p,+)\) 的群同态。若 \(c_\sigma=0\)，则 \(\sigma\) 固定生成元 \(\alpha\)，因而是恒等映射，所以这个同态是单射。目标群阶为素数，只有平凡子群和整个群，于是 \([L:K]\) 只能为 \(1\) 或 \(p\)。

次数为 \(1\) 时 \(\alpha\in K\)，所有根都在 \(K\) 中，且 \(a=\wp(\alpha)\)。反过来，若 \(a=\wp(b)\)、\(b\in K\)，全部根就是 \(b+c\)、\(c\in\mathbb F_p\)，故多项式完全分裂。次数为 \(p\) 时，\(\alpha\) 的最小多项式次数为 \(p\)，且整除首一 \(p\) 次多项式 \(X^p-X-a\)，因而二者相等，给定多项式不可约。此时上述单同态的像是整个 \((\mathbb F_p,+)\)，所以 Galois 群为循环群，并且每个平移 \(\alpha\mapsto\alpha+c\) 都对应一个 \(K\) 自同构。
\end{proof}

\subsection{素数次循环扩张的分类}
\label{b1:subsec:2004:02}

\term[Artin-Schreier-dingli]{Artin--Schreier 定理}[Artin--Schreier theorem]说明上一小节的方程穷尽了特征 \(p\) 中的循环 \(p\) 次扩张，并给出避免重复的参数空间。

\begin{theorem}[Artin--Schreier]\label{b1:thm:artin-schreier-classification}
设 \(K\) 为特征 \(p>0\) 的域，并固定 \(K\) 的代数闭包 \(\overline K\)，记 \(\wp(b)=b^p-b\)（\(b\in K\)）。每个包含于 \(\overline K\) 的 \(p\) 次循环扩张 \(L/K\) 都形如
\[
L=K(\alpha),\qquad\alpha^p-\alpha=a\in K\setminus\wp(K).
\]
这些位于固定代数闭包中的子扩张与 \(\mathbb F_p\) 向量空间 \(K/\wp(K)\) 的一维子空间一一对应。具体地，对 \(a,a'\in K\setminus\wp(K)\)，分别在 \(\overline K\) 中选取满足 \(\alpha^p-\alpha=a\)、\(\beta^p-\beta=a'\) 的根，则 \(K(\alpha)=K(\beta)\) 当且仅当
\[
a'=ca+\wp(b)\qquad\text{对某个 }c\in\mathbb F_p^\times,\ b\in K.
\]
\end{theorem}

在\cref{b1:thm:kummer-cyclic}的构造中，为得到 \(\sigma(\alpha)=\zeta\alpha\)，把 \(\dfrac{\alpha}{\sigma(\alpha)}=\zeta^{-1}\) 交给乘法形式的 Hilbert 第 90 定理。这里设 \(\sigma\) 是 \(p\) 次循环扩张 \(L/K\) 的生成自同构，改为要求 \(\sigma(\alpha)=\alpha+1\)，也就是 \(\alpha-\sigma(\alpha)=-1\)。这样 \(\alpha^p-\alpha\) 会被 \(\sigma\) 固定，而 \(-1\) 的迹为零，加法形式的 Hilbert 第 90 定理便能给出所需的 \(\alpha\)。

\begin{proof}
设 \(L/K\) 是 \(p\) 次循环扩张，生成自同构为 \(\sigma\)。由于扩张次数和基域特征都为 \(p\)，有 \(\operatorname{Tr}_{L/K}(-1)=-p\cdot1_K=0\)。\cref{b1:thm:hilbert90-additive}不要求扩张次数在 \(K\) 中可逆，其证明没有除以 \(p\)，所以这里仍可应用。它给出 \(\alpha\in L\) 满足 \(\alpha-\sigma(\alpha)=-1\)，即 \(\sigma(\alpha)=\alpha+1\)。于是
\[
\sigma(\alpha^p-\alpha)
=(\alpha+1)^p-(\alpha+1)
=\alpha^p+1-\alpha-1
=\alpha^p-\alpha.
\]
因为 \(\sigma\) 生成整个 Galois 群，这个元素属于固定域 \(K\)，记为 \(a\)。又因 \(\sigma(\alpha)\neq\alpha\)，有 \(\alpha\notin K\)。\cref{b1:prop:artin-schreier-polynomial}因而给出 \(a\notin\wp(K)\) 和 \([K(\alpha):K]=p\)。结合 \(K(\alpha)\subseteq L\) 及 \([L:K]=p\)，得到 \(K(\alpha)=L\)。

若 \(a'=ca+\wp(b)\)，其中 \(c\in\mathbb F_p^\times\)、\(b\in K\)，令 \(\gamma=c\alpha+b\)。由 \(c^p=c\) 得
\[
\gamma^p-\gamma=c(\alpha^p-\alpha)+(b^p-b)=a'.
\]
又因 \(\alpha=c^{-1}(\gamma-b)\)，有 \(K(\gamma)=K(\alpha)\)。两根 \(\beta,\gamma\) 相差 \(\mathbb F_p\) 中的元素，所以 \(K(\beta)=K(\gamma)=K(\alpha)\)。

反过来，设 \(K(\alpha)=K(\beta)=L\)。由\cref{b1:prop:artin-schreier-polynomial}，映射 \(\operatorname{Gal}(L/K)\to(\mathbb F_p,+)\)、\(\tau\mapsto\tau(\alpha)-\alpha\) 是同构。因此可选取 \(\sigma\) 使 \(\sigma(\alpha)=\alpha+1\)；它对应非零元素 \(1\)，故是 Galois 群的生成元。因为 \(a'\in K\)，有
\[
\sigma(\beta)^p-\sigma(\beta)=\sigma(\beta^p-\beta)=a'.
\]
这说明 \(\beta\) 与 \(\sigma(\beta)\) 是同一个 Artin--Schreier 方程的根，于是 \(\sigma(\beta)-\beta=c\in\mathbb F_p\)。若 \(c=0\)，则 \(\beta\) 被生成元 \(\sigma\) 固定，从而 \(\beta\in K\)，与 \(a'\notin\wp(K)\) 矛盾。因此 \(c\neq0\)，且
\[
\sigma(\beta-c\alpha)=\beta+c-c(\alpha+1)=\beta-c\alpha.
\]
故 \(\beta-c\alpha\) 等于某个 \(b\in K\)。代入 \(\beta=c\alpha+b\) 就得到 \(a'=ca+\wp(b)\)。

在商空间 \(K/\wp(K)\) 中，这个参数条件等价于 \([a']=c[a]\)，其中 \(c\in\mathbb F_p^\times\)。因此同一个扩张对应同一个一维子空间中的所有非零向量，而 \(\wp(b)\) 只改变向量的代表。每个非零类 \([a]\) 又由前一命题给出一个 \(p\) 次循环扩张，所以上述等价条件恰好给出所述的一一对应。
\end{proof}

例如在 \(K=\mathbb F_p(t)\) 中，\(X^p-X-t\) 不可约。由\cref{b1:prop:artin-schreier-polynomial}，只需证明 \(t\notin\wp(K)\)。若 \(b=0\)，显然 \(b^p-b\neq t\)；否则把 \(b\) 写成既约分式 \(\dfrac{u(t)}{v(t)}\)，其中 \(u,v\in\mathbb F_p[t]\) 都非零，令 \(m=\deg u-\deg v\)。此时
\[
b^p-b=\frac{u^p-uv^{p-1}}{v^p}.
\]
若 \(m>0\)，则 \(p\deg u>\deg u+(p-1)\deg v\)，所以分子的两项最高次项不会相消，分子次数为 \(p\deg u\)，分母次数为 \(p\deg v\)，次数差为 \(pm\)。约去公因子会使分子、分母次数减少同一个数，不改变这个差；而 \(t\) 的次数差为 \(1\)，不可能等于 \(pm\)。若 \(m\le0\)，分子次数不超过分母次数，或者分子为零，也不可能表示 \(t\)。这就证明了所需结论。通常把非零有理函数 \(\dfrac{u(t)}{v(t)}\) 的 \(\max(\deg u-\deg v,0)\) 称为它在无穷远处的极点阶；上述计算也说明，\(b\) 有正极点阶 \(m\) 时，\(b^p-b\) 的极点阶恰为 \(pm\)。

\subsection{与 Kummer 理论的比较}
\label{b1:subsec:2004:03}

设 \(\ell\) 为不同于 \(p\) 的素数，且 \(K\) 包含全部 \(\ell\) 次单位根。Kummer 理论使用乘法商群 \(K^\times/K^{\times\ell}\)，把它看成 \(\mathbb F_\ell\) 向量空间时，标量作用由取幂给出。方程为 \(X^\ell-a\)，根之间相差 \(\mu_\ell\) 中的乘法因子；选取本原 \(\ell\) 次单位根 \(\zeta\)，生成自同构可写为 \(\sigma(\alpha)=\zeta\alpha\)，其构造来自乘法形式的 Hilbert 第 90 定理。在特征 \(p\) 中，Artin--Schreier 理论使用加法商空间 \(K/\wp(K)\)，它按通常的加法与素域标量作用成为 \(\mathbb F_p\) 向量空间。方程为 \(X^p-X-a\)，根之间相差一个加法常数 \(c\in\mathbb F_p\)，生成自同构可写为 \(\sigma(\alpha)=\alpha+1\)，其构造来自加法形式的 Hilbert 第 90 定理。前者用 \(\dfrac{b}{\sigma(b)}\) 描述范数为一的元素，后者用 \(b-\sigma(b)\) 描述迹为零的元素；两者的非平凡素数次循环扩张都由相应商空间的一维子空间分类。

\ProseParagraphBreak
这也解释了为什么不能把特征零的根式可解性判据原样搬到正特征：特征 \(p\) 中的可分循环 \(p\) 次步骤由 Artin--Schreier 方程描述，\(p\) 次开方却是纯不可分步骤。二者的次数可以相同，自同构行为则完全不同。下一章先在特征零下给出根式可解性的完整等价定理。


\begin{chaptersummary}
Hilbert 第 90 定理将范数为一与共轭比值相联系，将迹为零与共轭差相联系。范数和迹均按全部自同构计算，相同的取值也保留出现次数。乘法证明构造一个非零的加权嵌入和，加法证明比较两个线性空间的维数；两个证明均不需要除以扩张次数。

\ProseParagraphBreak
当基域 \(K\) 的特征不整除 \(n\)，且 \(K\) 含 \(\mu_n\) 时，Kummer 配对把指数整除 \(n\) 的有限交换 Galois 扩张与有限根式数据相配。特征 \(p\) 中则由 \(X^p-X-a\) 代替 \(p\) 次开方。两类素数次扩张的分类参数都是相应商空间中的直线。这个“直线”的条件排除了因改变生成元而重复计算同一个扩张。
\end{chaptersummary}

\BookExercises

\begin{exercise}\label{b1:ex:20-1}
设 \(L=\mathbb F_{q^n}\)、\(K=\mathbb F_q\)。把范数视为 \(N_{L/K}:L^\times\rightarrow K^\times\)，求其核的阶；再证明映射 \(L^\times\rightarrow L^\times\)、\(b\mapsto b/b^q\) 的像恰为此核，并求每个像的原像数。
\end{exercise}
\begin{solution}
\cref{b1:thm:finite-field-automorphisms}给出 Frobenius \(b\mapsto b^q\) 生成 Galois 群，故\eqref{b1:eq:cyclic-trace-norm}中的范数为 \(b^{1+q+\cdots+q^{n-1}}\)。\cref{b1:cor:finite-field-cyclic}说明乘法群为 \(q^n-1\) 阶循环群，所以该幂映射的像阶为 \(q-1\)，核阶为 \((q^n-1)/(q-1)\)。映射 \(b\mapsto b^{1-q}\) 的核为 \(K^\times\)，也有 \(q-1\) 个元素；相消计算使其像包含于范数核，而像的阶已经等于范数核的阶，因而相等。每个非空纤维为核的一个陪集，所以有 \(q-1\) 个元素。
\end{solution}

\begin{exercise}\label{b1:ex:20-2}
设 \(d\in\mathbb Q^\times\) 不是平方。参数化全部有理数解 \((x,y)\) 满足 \(x^2-dy^2=1\)，并解释它们与二次扩张的范数有何关系。
\end{exercise}
\begin{solution}
左边为 \(N_{\mathbb Q(\sqrt d)/\mathbb Q}(x+y\sqrt d)\)。除去解 \((-1,0)\)，可令 \(t=y/(x+1)\)。代入 \((x-1)(x+1)=dy^2\) 并除以 \(x+1\)，得 \(x-1=dt^2(x+1)\)。因为 \(d\) 不是有理平方，\(1-dt^2\neq0\)，故
\[
x=\frac{1+dt^2}{1-dt^2},\qquad y=\frac{2t}{1-dt^2}\quad(t\in\mathbb Q).
\]
反向代入给出恒等式 \(x^2-dy^2=1\)。在扩张域中，这一参数化正是 \((1+t\sqrt d)/(1-t\sqrt d)\) 的有理化结果。
\end{solution}

\begin{exercise}\label{b1:ex:20-3}
令 \(L=\mathbb F_2(\theta)\)，其中 \(\theta^2+\theta+1=0\)。分别求出对 \(\mathbb F_2\) 的迹为零的元素集合、范数为一的元素集合，再求两集合的交，并解 \(b-b^2=1\)。
\end{exercise}
\begin{solution}
域中四个元素为 \(0,1,\theta,\theta+1\)，Frobenius 为 \(b\mapsto b^2\)。迹分别为 \(0,0,1,1\)，所以迹零集合为 \(\{0,1\}\)。每个非零元素的立方为 \(1\)，故范数一集合为 \(\{1,\theta,\theta+1\}\)，两集合的交为 \(\{1\}\)。方程 \(b-b^2=1\) 在特征 \(2\) 中就是 \(b+b^2=1\)，解为 \(\theta,\theta+1\)。这两个解相差素域元素，与\cref{b1:thm:hilbert90-additive}的陪集描述一致。
\end{solution}

\begin{exercise}\label{b1:ex:20-4}
设 \(L/K\) 为 \(n\) 次循环扩张，\(\sigma\) 生成其 Galois 群，\(b\in L^\times\)，并令 \(a=b/\sigma(b)\)。对正整数 \(r\)，证明
\[
\frac b{\sigma^r(b)}=\prod_{j=0}^{r-1}\sigma^j(a).
\]
写出加法版本。若 \(\gcd(r,n)=1\)，说明改用生成元 \(\sigma^r\) 为何不会改变两个定理所描述的核。
\end{exercise}
\begin{solution}
把各 \(\sigma^j(a)=\sigma^j(b)/\sigma^{j+1}(b)\) 相乘，分子分母相消即可。加法版本是 \(b-\sigma^r(b)=\sum_{j=0}^{r-1}\sigma^j\bigl(b-\sigma(b)\bigr)\)。若 \(\gcd(r,n)=1\)，则 \((\sigma^r)^i\)（\(0\leq i<n\)）仍遍历全部自同构，按这 \(n\) 项计算的迹与范数不变。由 Hilbert 第 90 定理的两个版本，两种生成元分别给出的像都等于相同的迹核、范数核。
\end{solution}

\begin{exercise}\label{b1:ex:20-5}
设 \(L=\mathbb Q(\sqrt2,\sqrt3)\)，\(\sigma\) 固定 \(\sqrt3\) 而把 \(\sqrt2\) 变号。构造范数 \(N_{L/\mathbb Q}\) 为 \(1\) 的元素，却不能写成 \(b/\sigma(b)\)。
\end{exercise}
\begin{solution}
取 \(a=(2+\sqrt3)/(2-\sqrt3)\)。它在 \(\mathbb Q(\sqrt3)/\mathbb Q\) 中的范数为 \(1\)，故由\cref{b1:thm:trace-norm-tower}，\(N_{L/\mathbb Q}(a)=1\)。若 \(a=b/\sigma(b)\)，则 \(a\cdot\sigma(a)=1\)。但 \(\sigma(a)=a\)，而 \(a\neq\pm1\)，故 \(a^2\neq1\)，矛盾。范数到 \(\mathbb Q\) 涉及四个自同构，单个二阶 \(\sigma\) 不能生成全部群，因而不满足乘法定理的假设。
\end{solution}

\begin{exercise}\label{b1:ex:20-6}
设 \(\ell\) 为不同于域 \(K\) 特征的素数，且 \(\mu_\ell\subseteq K\)。证明对任意 \(a\in K^\times\)，多项式 \(X^\ell-a\) 要么在 \(K\) 上完全分裂，要么不可约。

再设 \(L/K\) 为 \(d\) 次循环扩张，\(\sigma\) 生成其 Galois 群，\(K\) 含本原 \(d\) 次单位根 \(\zeta\)。给定一组 \(K\) 基 \(e_1,\ldots,e_d\)，令
\[
T_\zeta(x)=\sum_{i=0}^{d-1}\zeta^{-i}\sigma^i(x)\qquad(x\in L).
\]
证明至少有一个 \(T_\zeta(e_j)\) 非零，并证明任取这样的非零值 \(\alpha\)，都有 \(L=K(\alpha)\) 且 \(\alpha^d\in K\)。
\end{exercise}
\begin{solution}
由\cref{b1:prop:single-radical-cyclic}，任一根生成的扩张次数整除 \(\ell\)，故为 \(1\) 或 \(\ell\)。前者中一根属于 \(K\)，其余根由单位根相乘也属于 \(K\)；后者中最小多项式就是整个多项式。

\(T_\zeta\) 是 \(K\) 线性映射。\(\sigma^0,\ldots,\sigma^{d-1}\) 是不同的 \(K\) 嵌入，\cref{b1:cor:embedding-independence}说明它们以非零系数 \(1,\zeta^{-1},\ldots,\zeta^{-(d-1)}\) 所成的和不可能是零映射。因此 \(T_\zeta\) 不会在全部基向量上取零。重新排列求和指标可得 \(\sigma(T_\zeta(x))=\zeta T_\zeta(x)\)。若 \(\alpha=T_\zeta(e_j)\ne0\)，则 \(\alpha^d\) 被 \(\sigma\) 固定，故属于 \(K\)；其共轭元 \(\zeta^i\alpha\)（\(0\leq i<d\)）两两不同，所以 \([K(\alpha):K]\geq d=[L:K]\)，从而 \(L=K(\alpha)\)。
\end{solution}

\begin{exercise}\label{b1:ex:20-7}
求 \(\mathbb Q(\sqrt2,\sqrt3,\sqrt5)/\mathbb Q\) 的次数、Galois 群及全部二次中间域。
\end{exercise}
\begin{solution}
若 \(2^a3^b5^c\) 是有理平方，其中 \(a,b,c\in\{0,1\}\)，把有理数写成互素整数之比并比较各素数指数，得到 \(a=b=c=0\)。因此三个平方类独立，生成的子群阶为 \(8\)。\cref{b1:thm:kummer-correspondence}给出次数 \(8\)、群 \(C_2^3\)。二次中间域对应于平方类空间的一维子空间；在 \(\mathbb F_2\) 上每条直线只有一个非零元素，所以恰为 \(\mathbb Q(\sqrt d)\)，其中 \(d=2,3,5,6,10,15,30\)。
\end{solution}

\begin{exercise}\label{b1:ex:20-8}
设 \(K\) 为域，\(n\geq1\)，\(\operatorname{char}K\nmid n\)，且 \(\mu_n\subseteq K\)。取 \(a\in K^\times\) 及 \(\alpha^n=a\)。证明 \([K(\alpha):K]\) 等于 \([a]\) 在 \(K^\times/K^{\times n}\) 中的阶。再取 \(K=\mathbb Q(i)\)、\(n=4\)，分别以 \(a=2\) 和 \(a=4\) 计算扩张次数。
\end{exercise}
\begin{solution}
记 \(L=K(\alpha)\)、\(d=[L:K]\)，并令 \(e\) 为 \([a]\) 的阶。由\cref{b1:prop:single-radical-cyclic}，\(L/K\) 循环 Galois，且 \(\operatorname{Gal}(L/K)\) 嵌入 \(\mu_n\)。每个 \(\sigma\) 所对应的比值 \(\sigma(\alpha)/\alpha\) 都属于这个阶为 \(d\) 的子群，故 \(\sigma(\alpha^d)=\alpha^d\)。因此 \(\alpha^d\in K^\times\)，于是 \(a^d=(\alpha^d)^n\in K^{\times n}\)，得到 \(e\mid d\)。反过来，\(a^e=c^n\) 对某个 \(c\in K^\times\) 成立，因而 \((\alpha^e/c)^n=1\)。由 \(\mu_n\subseteq K\) 得 \(\alpha^e\in K\)，故 \(d\leq e\)。结合两式，\(d=e\)。

在 \(K=\mathbb Q(i)\) 中，\(4\notin K^{\times4}\)：若 \(c^4=4\)，则 \(c\in\{\pm\sqrt2,\pm i\sqrt2\}\)，这些元素均不在 \(K\) 中。因此 \([4]\) 的阶为 \(2\)，而 \([2]^2=[4]\ne[1]\)、\([2]^4=[16]=[1]\)，所以 \([2]\) 的阶为 \(4\)。相应地，四次方根 \(\alpha^4=2\) 生成四次扩张，\(\beta^4=4\) 生成二次扩张。合数指数下，根式的次数可以是真因子。
\end{solution}

\begin{exercise}\label{b1:ex:20-9}
设 \(K\) 为域，\(n\geq2\)，\(\operatorname{char}K\nmid n\)，且 \(\mu_n\subseteq K\)。在 \(K\) 的一个固定代数闭包中，对 \(K^\times/K^{\times n}\) 的任一有限子群 \(B\)，从每个类 \([a]\in B\) 中任选代表 \(a\in K^\times\) 和一个 \(n\) 次根，记 \(L_B=K(\sqrt[n]{a}\mid [a]\in B)\)。若 \(B_1,B_2\) 是两个这样的子群，证明 \(L_{B_1B_2}=L_{B_1}L_{B_2}\)，以及 \(L_{B_1\cap B_2}=L_{B_1}\cap L_{B_2}\)。
\end{exercise}
\begin{solution}
乘积子群由两个子群的生成元合在一起生成，加入它们的根得到的就是合成域，因此第一式成立。两域的交仍为有限 Galois 扩张，其群是其中任一 Galois 群的商，因而交换且指数整除 \(n\)。对 \(a\in K^\times\)，若其某个 \(n\) 次根分别属于两域，因为所有根只相差 \(K\) 中的单位根，任取一个根就同时属于两域。故交域对应的子群恰为 \(B_1\cap B_2\)，由对应的唯一性得到第二式。
\end{solution}

\begin{exercise}\label{b1:ex:20-10}
设 \(K\) 含本原三次单位根，并取 \(a\in K^\times\setminus K^{\times3}\)。对 \(\alpha^3=a\)，写出 \(K(\alpha)/K\) 的自同构，并证明 \(a^2b^3\)（\(b\in K^\times\)）给出同一个扩张。
\end{exercise}
\begin{solution}
\cref{b1:prop:single-radical-cyclic}及素数次数给出三个自同构 \(\alpha\mapsto\zeta_3^j\alpha\)。元素 \(b\alpha^2\) 的立方为 \(a^2b^3\)。反过来，由 \(\alpha=a^{-1}(\alpha^2)^2\)，得 \(K(b\alpha^2)=K(\alpha)\)。商群中 \([a]\) 和 \([a]^2\) 是同一直线上的两个不同非零元素，因此以单个非零类分类会重复计算这个域。
\end{solution}

\begin{exercise}\label{b1:ex:20-11}
在 \(K=\mathbb F_3(t)\) 上，比较参数 \(t\)、\(2t\)、\(t+t^3-t\) 给出的 Artin--Schreier 扩张，并说明它们非平凡。
\end{exercise}
\begin{solution}
第三个参数就是 \(t^3\)，与 \(t\) 的差为 \(\wp(t)\)；第二个参数是第一个的非零素域倍数。由\cref{b1:thm:artin-schreier-classification}，三者给出同一扩张。具体地，若 \(\alpha^3-\alpha=t\)，则 \(2\alpha\) 对应 \(2t\)，\(\alpha+t\) 对应 \(t^3\)。\(t\) 在无穷远处有一阶极点；若 \(\wp(b)\) 在无穷远处有极点，则其极点阶是 \(3\) 的倍数。故 \(t\notin\wp(K)\)，扩张次数为 \(3\)。
\end{solution}

\begin{exercise}\label{b1:ex:20-12}
设 \(p>2\)、\(K=\mathbb F_p(t)\)。证明 \(t,t^2\) 的类在 \(K/\wp(K)\) 中线性无关，并求由 \(\alpha^p-\alpha=t\)、\(\beta^p-\beta=t^2\) 生成的合成域次数及 Galois 群。列出合成域中全部对 \(K\) 为 \(p\) 次的中间域，并说明为什么没有遗漏。
\end{exercise}
\begin{solution}
任何非零组合 \(ct+dt^2\) 的无穷远极点阶为 \(1\) 或 \(2\)，不被 \(p\) 整除；它不可能等于 \(\wp(b)\)，故两类独立。两个 \(p\) 次循环域因其参数直线不同而不同。把第二个 Artin--Schreier 多项式看作第一个域上的多项式，由\cref{b1:prop:artin-schreier-polynomial}，其根要么已在该域，要么产生 \(p\) 次扩张；前者会使两个原来的 \(p\) 次域相等，已被排除。因此合成域 \(M=K(\alpha,\beta)\) 的次数为 \(p^2\)。它是两个可分多项式的共同分裂域，限制到两域给出到 \(C_p\times C_p\) 的单射，两边阶相同，所以 Galois 群正是该直积。

若 \(c\in\mathbb F_p\)，则 \((\alpha+c\beta)^p-(\alpha+c\beta)=t+ct^2\)，其参数类与 \([t^2]\) 一起代表二维空间 \(\langle[t],[t^2]\rangle_{\mathbb F_p}\) 的全部 \(p+1\) 条直线。因此\cref{b1:thm:artin-schreier-classification}给出 \(p+1\) 个互异的 \(p\) 次中间域
\[
K(\beta),\qquad K(\alpha+c\beta)\quad(c\in\mathbb F_p).
\]
另一方面，\(\operatorname{Gal}(M/K)\cong C_p^2\) 恰有 \(p+1\) 个指数为 \(p\) 的子群；由 Galois 对应，没有其他 \(p\) 次中间域。
\end{solution}

\begin{exercise}\label{b1:ex:20-13}
令 \(K=\mathbb F_{p^r}\)。证明 \(\wp(K)=\ker\operatorname{Tr}_{K/\mathbb F_p}\)，并据此说明 \(K\) 在固定代数闭包中恰有一个 \(p\) 次循环扩张。
\end{exercise}
\begin{solution}
因 \(\ker\wp=\mathbb F_p\)，\(\wp(K)\) 有 \(p^{r-1}\) 个元素。迹的嵌入和公式给出 \(\operatorname{Tr}(b^p)=\operatorname{Tr}(b)\)，因为取 \(p\) 次幂只会循环置换各项，所以 \(\wp(K)\) 包含于迹核。由\cref{b1:cor:embedding-independence}，这些不同嵌入之和不是零映射，而目标空间 \(\mathbb F_p\) 为一维，故迹满射。迹核也有 \(p^{r-1}\) 个元素，两者相等。商 \(K/\wp(K)\) 因而是一维 \(\mathbb F_p\) 空间，只有一条一维子空间。应用\cref{b1:thm:artin-schreier-classification}即可。
\end{solution}

\begin{exercise}\label{b1:ex:20-14}
证明 \(X^p-X-1\) 在 \(\mathbb F_p[X]\) 中不可约，并求其根 \(\alpha\) 到 \(\mathbb F_p\) 的迹和范数。
\end{exercise}
\begin{solution}
\(\wp(\mathbb F_p)=0\)，故\cref{b1:prop:artin-schreier-polynomial}给出不可约性。全部根为 \(\alpha+c\)、\(c\in\mathbb F_p\)，其和为 \(\sum c\)。该和在 \(p>2\) 时为 \(0\)，在 \(p=2\) 时为 \(1\)。范数是根的乘积，由首一多项式常数项公式等于 \((-1)^p(-1)=1\) 于 \(\mathbb F_p\)。迹在奇素数情形为零，不妨碍扩张可分或迹映射满射。
\end{solution}

\begin{exercise}\label{b1:ex:20-15}
设 \(K=\mathbb F_p(t)\)、\(L=K(u)\)，其中 \(u^p=t\)。证明 \([L:K]=p\) 且 \(\operatorname{Aut}_K(L)=1\)，并与参数 \(t\) 的 Artin--Schreier 扩张比较。
\end{exercise}
\begin{solution}
在 \(\mathbb F_p[t][X]\) 中，\(X^p-t\) 对素元 \(t\) 满足\cref{b1:thm:eisenstein}的条件：非首项系数均被 \(t\) 整除，常数项不被 \(t^2\) 整除。因此它在 \(K[X]\) 中不可约，次数为 \(p\)。任一自同构必须把 \(u\) 送到 \(X^p-t=(X-u)^p\) 的根，唯一可能是 \(u\)，所以只能是恒等映射。相比之下，\(X^p-X-t\) 的导数为 \(-1\)，而 \(t\) 的无穷远极点阶为 \(1\)，不能是 \(\wp(b)\) 的极点阶，所以\cref{b1:prop:artin-schreier-polynomial}给出其不可约性；其根生成同次数的可分正规扩张，自同构群为 \(C_p\)。
\end{solution}
